\section{Why Fuzz Cryptographic Code}


\subsection{The Problem with Test Vectors}

Cryptographic implementations are typically tested against known-answer tests (KATs):
standardized input/output pairs from the specification. KATs verify that the code
produces correct output for specific inputs. They do not verify:

\begin{itemize}[nosep]
  \item Behavior at boundary values (zero, one, $q - 1$, $q$, $q + 1$).
  \item Behavior with malformed inputs (wrong length, invalid encoding).
  \item Absence of panics or undefined behavior.
  \item Consistency: $\text{decrypt}(\text{encrypt}(m)) = m$ for all $m$.
  \item Algebraic properties: associativity, commutativity, distributivity.
\end{itemize}

\subsection{Property-Based Testing}

Property-based testing (PBT) inverts the test structure. Instead of specifying inputs
and expected outputs, you specify \emph{properties} that must hold for all inputs. The
framework generates random inputs and checks the properties.

\begin{definition}[Property]
A property is a universally quantified assertion: $\forall x \in D.\; P(x)$, where $D$
is the input domain and $P$ is a predicate.
\end{definition}

Examples of cryptographic properties:
\begin{itemize}[nosep]
  \item \textbf{Round-trip}: $\text{decrypt}(k, \text{encrypt}(k, m)) = m$
  \item \textbf{Determinism}: $\text{hash}(x) = \text{hash}(x)$
  \item \textbf{Length preservation}: $|\text{NTT}(a)| = |a|$
  \item \textbf{Commutativity}: $\text{NTT}(a \cdot b) = \text{NTT}(a) \otimes \text{NTT}(b)$
  \item \textbf{Idempotence}: $\text{NTT}(\text{INTT}(\text{NTT}(a))) = \text{NTT}(a)$
\end{itemize}

\subsection{Fuzzing vs.\ Property-Based Testing}

\begin{table}[H]
\centering
\small
\begin{tabular}{@{}lll@{}}
\toprule
 & \textbf{Property-Based Testing} & \textbf{Fuzzing} \\
\midrule
Input generation & Type-aware random & Mutational, coverage-guided \\
Oracle & Explicit property & Crash / sanitizer \\
Shrinking & Yes (minimal failing input) & Corpus minimization \\
Coverage feedback & Usually no & Yes \\
Typical runs & 1,000--100,000 & Millions--billions \\
\bottomrule
\end{tabular}
\caption{PBT vs.\ fuzzing. Modern tools blend both approaches.}
\end{table}

Go's \texttt{testing.F} provides coverage-guided fuzzing with type-aware seed corpora.
Foundry provides both random (fuzz) and coverage-guided (invariant) modes. We use all
of them.

%% ============================================================
